\section{Motivation} 
\label{sec:introduction:motivation}

\begin{highlightbox}
In this seemingly certain world, uncertainty lingers like a shadow, yet from the uncertain, we often find certainty.       
\end{highlightbox}  

The paradoxical relationship between certainty and uncertainty appears across disciplines, demonstrating how truth often arises from navigating ambiguity. Quantum mechanics redefines predictability under randomness in the subatomic world, where Newtonian deterministic physics collapses. This uncertainty led to groundbreaking laws (e.g. Heisenberg's Uncertainty Principle), which now underpin modern technologies like quantum computing and all forms of spectroscopy. In medical diagnosis, for example, a patient may present vague symptoms, such as fatigue or pain, that could indicate anything from stress to cancer. Doctors converge on the cause through iterative hypothesis-testing approaches. Climate science also epitomizes how we manage uncertainty. Climate models incorporate a wide range of sources of variability and uncertainty. A common way to address this uncertainty is to combine the results of different simulation models. By running multi-model ensembles, scientists can identify consistent patterns across the uncertainty.

From quantum randomness to medical diagnosis to climate models, uncertainty does not destroy deterministic systems, but instead provides access to identifying them. This perspective supports the use of probabilistic inference methods, which offer a rigorous framework for representing uncertainty and reasoning from uncertainty. In statistics, modeling the uncertainty of variables from given data samples using a probability distribution rather than a point estimate is known as inference \footnote{We note that "inference" is an overloaded term in this context. We distinguish between: (i) Bayesian Inference, the statistical goal of recovering the posterior distribution; (ii) Variational Inference, the mechanism to achieve Bayesian inference; (iii) Inference: the process that maps data to latent space in generative models, which is in contrast to generation; and (iv) Computational Inference, the deep learning engineering practice of executing the generation pass (for generative models meaning from latent to data) of a trained network on test data.}. The estimated distribution facilitates uncertainty quantification and provides a measure of the reliability of the results, making it particularly well suited for real-world applications. 

Probabilistic inference grounded in Bayes' rule is known as {\bfseries\itshape Bayesian inference}. Bayes' rule represents unknown random variables through beliefs (the prior). The beliefs are updated using observed data (likelihood), yielding the posterior distribution. A simplified formulation is:
\begin{equation*}
    \label{eq:bayes_sim}
    \text{posterior} = \text{prior}\times \text{likelihood} \;/\; \text{total evidence (or marginal likelihood)}.
\end{equation*}
For example, we can illustrate Bayes' rule using the medical diagnosis scenario: the objective is to obtain the probability of a disease given a positive test result:
\begin{align*}
& P(\text{Disease} | \text{Positive}) = \frac{P(\text{Positive} | \text{Disease}) \times P(\text{Disease})}{P(\text{Positive})},\\
& P(\text{Positive}) = P(\text{Positive} | \text{Disease}) \times P(\text{Disease}) + P(\text{Positive} | \overline{\text{Disease}}) \times P(\overline{\text{Disease}}),
\end{align*}
where $\overline{\text{Disease}}$ denotes "No Disease".

In a more general Bayesian context, inference is the process of obtaining the posterior distribution of target variables from observed data. This perspective extends naturally to modern machine learning, where unknown values, structures, and policies can all be understood as random variables and probability distributions \citep{pml1Book}.

Especially in \acfp{DGM}, Bayesian inference plays an essential role \citep{pml2Book}. By introducing random variables to represent complex data in latent space, \acp{DGM} enable data synthesis from latent variables (generation) based on recognition of the underlying latent representation from data (inference). However, directly inferring the posterior distribution over the latent variables using Bayes' rule is generally intractable, since all possible values of the latent variables are rarely obtainable in complex and high-dimensional models. Addressing this challenge is essential to fully realize the potential of \acp{DGM}. \acfp{VAE} \citep{Kingma_2019} tackle this problem by introducing an inference network as a tractable approximation to the posterior. \acfp{NF} \citep{rezende2016} construct bijective transformations between observed data and latent variables, achieving exact computation of posteriors. Diffusion models \citep{ho2020ddpm} sidestep direct calculation of the posterior distribution over latent variables by designing it as a simple Gaussian, and frame both inference and generation as an iterative denoising process from noise through learned transformations.

In summary, \acp{DGM} are inherently probabilistic; through probabilistic modeling and inference, they are powerful tools for extracting latent information and representing uncertainty in complex data. These capabilities have been demonstrated in numerous applications involving images, text, and speech. The same principles extend naturally to scientific and engineering domains. One prominent question is: how can \acp{DGM} support the identification and modeling of complex, stochastic, nonlinear dynamical systems? In control theory, the process of building mathematical models of dynamical systems from observed input-output data is called {\bfseries\itshape system identification}. Interpreting system identification in probabilistic terms, {\bfseries\itshape dynamics modeling} means inferring the posterior of the next states based on the previous states and system inputs, while {\bfseries\itshape parameter identification} means estimating the posterior over system parameters given observations. Identifying systems probabilistically enables uncertainty quantification, which is crucial for reliable and robust decision-making in real-world applications.

Classical methods for system identification rely on linearized models or prespecified nonlinear structures. Model parameters are then estimated from observed data using estimation techniques. Subsequently, a validation process is employed to assess and select the most appropriate model structure. Although these methods offer strong physical interpretability and are often data-efficient, they exhibit several limitations: parameter estimation can be computationally inefficient with large amounts of data; results are sensitive to the choice of model structures and depend on expert knowledge; uncertainty is difficult to handle and quantify; and they are generally inapplicable to partially observable systems and high-dimensional raw data. 

Alongside the demand for uncertainty quantification, deep learning methods are well suited for system identification because of their strength in representing information from complex data, their flexibility in function approximation, and their scalability to high-dimensional data. Furthermore, their compatibility with gradient-based optimization enables efficient navigation of high-dimensional landscapes, allowing optimizers to find solutions nearly equivalent to the global optimum. This provides a distinct advantage over classical methods, where low-dimensional parameter spaces are more prone to being trapped in poor local minima that fail to capture the underlying dynamics. Taking uncertainty quantification into account, a simple yet effective technique is Gaussian Process Regression \citep{sarkka2019gp}, which involves assigning a Gaussian process prior to the latent regressor and performing statistical conditioning based on the available observations. This method has gained widespread adoption due to its nonparametric character and adaptability. Nonetheless, its significant computational and memory requirements limit its feasibility when dealing with large datasets. Another straightforward approach to combining the advantages of deep learning methods and uncertainty modeling is the \acf{BNN} \citep{Goan_2020}. Unlike standard neural networks that produce point estimates, \acp{BNN} assign probability distributions over the network parameters, thereby inducing distributions for outputs. However, \acp{BNN} show limitations in expressiveness because the approximate inference may not reliably capture true posterior behavior \citep{foong2020bnn}. Therefore, the adaptation and application of \acp{DGM} are essential for achieving rich representational capacity, flexible modeling frameworks, and the ability to effectively handle large-scale data.

Compared with classical Bayesian inference methods, \acp{DGM} provide a powerful route to amortized Bayesian inference. Instead of solving a separate inverse problem for each new observation, a single conditional generative model is optimized over many simulated or measured samples and then reused to infer posterior distributions for unseen cases. This amortization is particularly attractive for practical engineering datasets, where repeated parameter estimation must be performed under noise, model uncertainty, and large numbers of operating conditions, for example, in end-of-line testing of large numbers of products. 

Different \acp{DGM} have been investigated in system identification applications. Here, we briefly evaluate three prominent frameworks: \acp{VAE}, \acp{NF}, and diffusion models. \acp{VAE} work by compressing complex data into a lower-dimensional latent space and then learning to reconstruct original data from compressed latent variables. This makes \acp{VAE} flexible enough to work with time-series modeling architectures such as \acf{RNN} \citep{Girin_2021} and \acf{NARX} \citep{PANIAGUA2024200344}, which is beneficial for uncovering time-dependent features in data. However, \acp{VAE} suffer from a well-known problem of variance shrinkage \citep{asperti2020vae} and tend to converge to a degenerate local optimum, which is also known as posterior collapse \citep{wang2023vae}. Unlike \acp{VAE}, \acp{NF} take a different approach by using a series of invertible transformations to turn a simple probability distribution into a complex one that matches the data. The requirement of invertibility restricts their choices of neural network architectures, which limits their direct applicability to time series modeling. On the other hand, \acp{NF} provide exact and expressive density estimation capabilities. This means they can precisely calculate the mathematical probability of a specific outcome, making them exceptionally strong at quantifying uncertainty and enhancing the accuracy of probabilistic forecasts. Diffusion models have recently gained popularity by learning to generate high-quality and coherent data through a process of gradually reversing noise. They relax architectural constraints and show strong predictive generation capabilities. Nevertheless, their iterative data generation procedure results in substantial computational overhead. Balancing expressiveness, computational efficiency, and sample quality, {\bfseries\itshape normalizing flows} are chosen in this work as the generative backbone of probabilistic system identification.

While parameter identification focuses on determining the time-invariant system parameters, {\bfseries\itshape state estimation} is the process of inferring hidden, time-varying state variables from a history of noisy and potentially partial observations and known control inputs. Because of the noise in the sensors (measurement noise) and unmodeled disturbances in the system itself (process noise), state estimation is naturally probabilistic. Thus, estimating the target states means computing their posterior distribution. The traditional method for state estimation is the Luenberger State Observer, which assumes the system is deterministic and both the system dynamics and measurement models are linear and time-invariant. Taking stochasticity into account, the Kalman filter is the most widely applied method. Beyond system linearity, a Kalman filter assumes the process and measurement noises are additive zero-mean Gaussian noise. To handle nonlinear systems, the Extended Kalman filter linearizes the nonlinear functions around the state estimate. However, the linearization introduces errors, leading to incorrect estimates or even divergence. An improved modification is the Unscented Kalman filter, which carefully selects the means and variances and propagates them through the nonlinear functions. However, it is still only compatible with Gaussian noise.

By relaxing the requirement for rigid physical modeling and avoiding restrictive assumptions about uncertainty, while maintaining high expressiveness, modeling flexibility, and scalability to large datasets, the selected \ac{NF} architecture establishes a foundation for both probabilistic system identification and state estimation. From a data-centric perspective, these tasks represent different mapping challenges: parameter identification maps time series to a set of static data points; state estimation involves inferring data points or reconstructing temporal trajectories from noisy observations. To address these challenges, it is necessary to extend the standard normalizing flows, given the limited research currently available in this area. Bridging this gap not only broadens the applicability of \acp{NF} but also contributes to advancing Bayesian inference techniques for complex, time-varying systems.

Having explored system identification and state estimation through Bayesian inference, we can advance toward {\bfseries\itshape optimal control}. By leveraging data as the primary source of insight into system behavior, we enable control strategies that are more adaptive, scalable, and robust. This is especially advantageous in situations where analytical models are incomplete, uncertain, or impractical to derive. One direct application is model-free control, which obtains control inputs directly from observations and reference trajectories without requiring an explicit system model. This architecture is especially effective for real-time path planning and tracking. By bypassing complex analytical derivations, it allows for the rapid evaluation of numerous candidate trajectories, enabling the system to select the best path under time-sensitive constraints. In recent literature, diffusion models have attracted widespread attention for representing control policies due to their ability to generate high-quality trajectories under complex observations. However, while diffusion models excel in sample quality, normalizing flows provide more efficient sampling and a more rigorous probabilistic formulation. Despite these advantages, the use of \acp{NF} as a generative backbone for control remains largely underexplored. Building on model-free principles, \ac{RL} extends this paradigm by optimizing these stochastic policies through continuous interaction with the environment. \ac{RL} agents learn to maximize long-term rewards while adapting to system uncertainties, effectively enhancing the capabilities of traditional \ac{MPC}. Integrating \ac{RL} with \acp{NF} offers a powerful framework for achieving uncertainty-aware, robust, and self-improving control in dynamic environments.

In summary, this work is motivated by four key considerations:
\begin{itemize}
    \item Bayesian inference provides a principled framework for modeling uncertainty and extracting knowledge under uncertainty, enabling system identification and state estimation based solely on observations. This is particularly valuable in scenarios where first-principles models are unavailable, incomplete, or prohibitively complex.
    \item \acp{DGM} enable amortized Bayesian inference by learning reusable posterior estimators from many data samples, making them well suited to large-scale parameter identification tasks. The differences among \acp{DGM} in these tasks have not been fully assessed, particularly in terms of posterior accuracy, uncertainty calibration, identifiability, and expressiveness for multimodality.
    \item While \acfp{NF} offer the expressiveness, exact density estimation, and scalability required for complex dynamical systems, their application to system identification and state estimation remains significantly underexplored. Moreover, current \ac{NF} architectures are often not natively designed for temporal dependencies; therefore, specific adaptations to handle sequential data and state trajectories for these tasks need to be explored.
    \item Integrating \acfp{NF} in control and reinforcement learning frameworks enables adaptive, uncertainty-aware decision-making in complex, unpredictable environments. This approach has not been fully investigated and offers significant potential for advancing control strategies. 
\end{itemize}